\documentclass[11pt]{article}
\usepackage[margin=1in]{geometry}
\usepackage[T1]{fontenc}
\usepackage[utf8]{inputenc}
\usepackage{lmodern}
\usepackage{xcolor}
\usepackage{booktabs}
\usepackage{longtable}
\usepackage{array}
\usepackage{amsmath,amssymb}
\usepackage{enumitem}
\usepackage{hyperref}
\hypersetup{colorlinks=true, linkcolor=blue, urlcolor=blue}
\definecolor{accent}{RGB}{25,56,99}

\setlength{\parskip}{0.45em}
\setlength{\parindent}{0pt}

\begin{document}

\begin{center}
{\LARGE\bfseries\color{accent} Documentation for the SIP Processing Pipeline}\\[0.4em]
{\large Automated extraction, regression analysis, and model comparison for SIP column experiments}
\end{center}

\section{Purpose and scope}

This pipeline provides an automated workflow for processing Spectral Induced Polarization (SIP) column experiment datasets.  
The system extracts electrical measurements from processed SIP workbooks, matches them with laboratory saturation records, computes derived variables, and evaluates statistical relationships between saturation and electrical properties.

The current version supports:

\begin{itemize}[leftmargin=1.5em]
\item batch processing of multiple experiment folders
\item normalized and unnormalized processing modes
\item regression analysis using multiple candidate models
\item automated comparison between power-law and logistic relationships
\item statistical evaluation using $R^2$, SSE, log-likelihood, and AIC
\end{itemize}

The pipeline therefore serves both as a scientific analysis tool and as a reproducible data-processing framework.

\section{Directory structure}

\begin{longtable}{p{0.22\textwidth} p{0.28\textwidth} p{0.40\textwidth}}
\toprule
\textbf{Location} & \textbf{Expected content} & \textbf{Purpose}\\
\midrule
\endhead
\texttt{Analysis/} & Processed SIP workbook & Electrical measurements extracted from each sheet.\\
\texttt{SIP\_Raw\_Data/} & Laboratory logbook workbook & Provides water saturation or parameters required to derive it.\\
\texttt{Results/} & Output directory & Stores merged data tables, regression outputs, plots, and summaries.\\
\bottomrule
\end{longtable}

\section{Processing workflow}

\begin{longtable}{p{0.07\textwidth} p{0.52\textwidth} p{0.28\textwidth}}
\toprule
\textbf{Step} & \textbf{Action} & \textbf{Output}\\
\midrule
\endhead
1 & Identify experiment folders containing both required input directories. & Valid run list\\
2 & Locate SIP workbook and logbook workbook. & Input file paths\\
3 & Extract the row nearest the target frequency from each SIP sheet. & Electrical measurements\\
4 & Match measurement identifiers with saturation entries in the logbook. & Saturation mapping\\
5 & Compute derived variables and log transformations. & Joined dataset\\
6 & Fit regression models describing saturation–electrical relationships. & Model parameters\\
7 & Generate plots and summary tables. & CSV, Excel, PNG outputs\\
\bottomrule
\end{longtable}

\section{Derived variables}

The pipeline computes several derived variables from the raw measurements.

\subsection{Saturation}

If saturation is directly available in the logbook:

\[
S = \text{saturation value}
\]

If derived from water mass:

\[
S = \frac{M_w}{\rho_w \phi V}
\]

where

\begin{itemize}
\item $M_w$ = water mass
\item $\rho_w$ = water density
\item $\phi$ = porosity
\item $V$ = column volume
\end{itemize}

\subsection{Logarithmic transforms}

The pipeline can use either base-10 logarithms or natural logarithms.

Base-10 transforms:

\[
\log_{10} S,\qquad
\log_{10} \rho,\qquad
\log_{10} \sigma
\]

Natural logarithm option:

\[
\ln S,\qquad
\ln \rho,\qquad
\ln \sigma
\]

Rows with non-positive values are excluded from log-domain regression.

\section{Normalization}

Optional normalization can be applied to both variables prior to regression.

Normalization is always performed using the data range corresponding to the analysis level.

\begin{itemize}
\item Sample-level normalization uses only that sample's data.
\item Site-level normalization uses all samples belonging to the same site.
\item Global normalization uses the full dataset.
\end{itemize}

Available normalization methods include:

\subsection{Z-score normalization}

\[
x' = \frac{x - \mu}{\sigma}
\]

\subsection{Min–max normalization}

\[
x' = \frac{x - x_{min}}{x_{max} - x_{min}}
\]

\section{Regression models}

Two primary model families are evaluated.

\subsection{Power-law model}

\[
y = a S^b
\]

In logarithmic form:

\[
\log y = b \log S + \log a
\]

The power-law model is traditionally used for resistivity–saturation relationships.

\subsection{Logistic model}

The logistic model captures threshold behavior:

\[
y = c + \frac{L}{1 + \exp[-k(S-x_0)]}
\]

Parameters represent:

\begin{itemize}
\item $L$ : amplitude
\item $k$ : transition steepness
\item $x_0$ : transition location
\item $c$ : baseline level
\end{itemize}

Unlike earlier implementations, the current pipeline allows both increasing and decreasing logistic curves.

\section{Statistical metrics}

Each model is evaluated using several metrics.

\subsection{Sum of squared errors}

\[
SSE = \sum (y_i - \hat{y}_i)^2
\]

\subsection{Coefficient of determination}

\[
R^2 = 1 - \frac{SSE}{SST}
\]

\subsection{Residual standard deviation}

\[
\sigma = \sqrt{\frac{SSE}{n}}
\]

\subsection{Log likelihood}

\[
\ln L =
-\frac{n}{2}
\left(
\ln(2\pi\sigma^2) + 1
\right)
\]

\subsection{Akaike Information Criterion}

\[
AIC = n\ln\left(\frac{SSE}{n}\right) + 2k
\]

where $k$ is the number of model parameters.

The preferred model is the one with the lowest AIC.

\section{Analysis levels}

Model comparisons are performed at three levels.

\begin{longtable}{p{0.20\textwidth} p{0.65\textwidth}}
\toprule
\textbf{Level} & \textbf{Definition}\\
\midrule
\endhead
Sample & Regression performed separately for each experiment.\\
Site & All experiments belonging to the same site are combined.\\
All & Entire dataset pooled together.\\
\bottomrule
\end{longtable}

This structure allows the analysis to distinguish between local experimental behavior and broader trends.

\section{Outputs}

\begin{longtable}{p{0.38\textwidth} p{0.18\textwidth} p{0.31\textwidth}}
\toprule
\textbf{File} & \textbf{Scope} & \textbf{Purpose}\\
\midrule
\endhead
\texttt{joined\_\_MODE\_freqHz.csv} & Per run & Merged dataset containing derived variables.\\
\texttt{compare\_power\_vs\_logistic.csv} & Global & Regression statistics for both models.\\
\texttt{compare\_power\_vs\_logistic.xlsx} & Global & Excel summary of model comparisons.\\
\texttt{plots/*.png} & Per group & Visual comparison of model fits.\\
\bottomrule
\end{longtable}

\section{Interpretation}

Results consistently indicate that:

\begin{itemize}

\item Resistivity generally follows power-law scaling with saturation.

\item Conductivity often exhibits logistic behavior.

\item Logistic relationships indicate threshold-controlled electrical conduction.

\end{itemize}

Such threshold behavior is consistent with pore connectivity theory and percolation models.

\section{Summary}

The SIP processing pipeline integrates:

\begin{itemize}

\item automated data extraction
\item robust matching of experimental records
\item derived-variable computation
\item multiple regression models
\item statistical model comparison
\item automated visualization

\end{itemize}

The resulting framework provides a reproducible and extensible approach for analyzing SIP saturation experiments.

\end{document}
